% !TeX program = xelatex
\documentclass[12pt, a4paper, titlepage]{article}
\usepackage{xeCJK}
\usepackage{amsmath, amssymb}
\usepackage{geometry}
\usepackage{longtable}
\usepackage{booktabs}
\usepackage{tabularx}
\usepackage{caption}
\usepackage{titlesec}
\usepackage{silence}
\WarningFilter{caption}{Unknown document class}
\WarningFilter{xeCJK}{Redefining CJKfamily}
\WarningFilter{relsize}{Failed to get list}

% -- Fonts --
\setmainfont{Times New Roman}
\setCJKmainfont{Hiragino Mincho ProN}

% ページ設定
\geometry{a4paper, margin=2.5cm}
%\onehalfspacing

% セクションタイトルのフォーマット
\titleformat{\section}{\large\bfseries}{\thesection}{1em}{}
\titleformat{\subsection}{\normalsize\bfseries}{\thesubsection}{1em}{}

\begin{document}

\begin{center}
{\LARGE CCD-Balanced-Force法 コード実装と検証のためのテスト仕様書}
\end{center}

\vspace{1em}

本稿は、気液二相流解析手法「CCD-Balanced-Force法」の論文から実際の計算コードを開発するにあたり、実装されたコードがバグを含まず、理論通りの精度と物理的振る舞いを実現しているかを検証するためのテスト仕様書である。

開発プロセスを4つのフェーズに分割し、各関数・モジュールのユニットテスト（単体テスト）から、ソルバ全体の結合テストへと段階的に進めるボトムアップ・アプローチを規定する。

\section*{フェーズ1: 基盤数学モジュールとデータ構造の検証（Lower Layer）}

流体計算の前提となる、格子データの扱いと数学的な微分・補間スキームが正しく実装されているかを検証する。

\subsection*{1.1 ゴーストセルと境界条件のテスト}

\textbf{検証対象:} データ配列の確保と、境界条件（ディリクレ・ノイマン）の適用関数。

\textbf{検証方法:} 領域全体に定数や線形関数を代入し、境界条件関数を呼び出した後、ゴーストセル（6次精度CCDの場合は境界外側3セル分）に正しい値が外挿・コピーされているかを確認する。

\textbf{合否判定:} 境界を跨ぐ差分計算を行っても、領域端のセルで配列外参照エラー（Segfault）が起きず、かつ解析解と一致すること。

\subsection*{1.2 CCD法による微係数算出の単体テスト（関連: 4.1節）}

\textbf{検証対象:} 勾配（$\nabla$）およびラプラシアン（$\nabla^2$）を算出するCCD関数。

\textbf{検証方法:} $f(x,y) = x^2 + y^2$ などの関数を用いて、微分を計算し、理論値と比較する。

\textbf{合否判定:} 数値解と解析解の誤差が許容範囲内に収まること。

\section*{フェーズ2: 多相流インターフェースモジュールの検証}

気液界面を表現するレベルセット関数と、それに伴う物性値の急激な変化を処理する関数のテストを行う。

\subsection*{2.1 結合レベルセット（CLS）と再初期化のテスト（関連: 6.1節）}

\textbf{検証対象:} VOFとLSを結合する関数、および再初期化（Re-initialization）ソルバ。

\textbf{検証方法:} 人工的にノイズを乗せた符号付距離関数や、勾配が崩れたステップ関数を入力し、再初期化計算を数ステップ実行する。

\textbf{合否判定:} 界面位置（$\psi=0.5$ または $\phi=0$）が初期位置から移動しないこと（変位量 $< 10^{-4}\Delta x$）。界面近傍で $|\nabla \phi| = 1.0 \pm 10^{-3}$ が達成されること。

\subsection*{2.2 Heaviside関数と物性値補間のテスト（関連: 10.1節）}

\textbf{検証対象:} スムーズ化Heaviside関数 $H_\epsilon(\phi)$ と、密度 $\rho$・粘度 $\mu$ のセルフェースへの補間関数。

\textbf{検証方法:} 水（$\rho=1000$）と空気（$\rho=1$）の密度比を設定し、界面を横切る1次元ライン上で補間値をプロットする。

\textbf{合否判定:} 質量保存に直結する密度は算術平均等で滑らかに繋がっていること。圧力ポアソン方程式に用いる密度の逆数（$1/\rho$）は調和平均として正しく実装されており、界面で非物理的なスパイク（オーバーシュート）が発生していないこと。

\subsection*{2.3 界面幾何量とCSFモデルの体積力変換テスト（関連: 5.2節, 8.1節）}

\textbf{検証対象:} 曲率 $\kappa$ の算出関数、およびそれを表面張力体積力 $\mathbf{F}_{sv} = \sigma \kappa \nabla H_\epsilon$ に変換する関数。

\textbf{検証方法:} 半径 $R$ の静止円形界面を入力。

\textbf{合否判定:} $\kappa$ が解析解 $1/R$ に対し $O(h^6)$ の精度を持つこと。領域全体での体積力の積分値 $\int \mathbf{F}_{sv} dV$ がゼロベクトルになること（運動量保存の確認）。

\section*{フェーズ3: NS方程式ソルバモジュールの検証}

流体計算の各項が、物理法則に従って正しく時間発展するかを検証する。

\subsection*{3.1 移流スキーム（DCCD）の質量保存テスト（関連: 6.2節）}

\textbf{検証対象:} 対流項 $\nabla \cdot (\mathbf{u} \phi)$ および $\nabla \cdot (\mathbf{u} \mathbf{u})$ を解くDCCD関数。

\textbf{検証方法:} Zalesak円盤テスト。規定の回転速度場 $\mathbf{u}$ を与え、スリット入り円盤を1回転させる。

\textbf{合否判定:} 初期形状に対する質量欠損率が $0.1\%$ 未満であり、スリットの角の丸まりが視覚的・定量的に閾値内であること。

\subsection*{3.2 粘性拡散スキームのテスト（関連: 7.1節）}

\textbf{検証対象:} 粘性項 $\nabla \cdot (\mu \nabla \mathbf{u})$ の陰解法または陽解法ソルバ。

\textbf{検証方法:} 平行平板間のポアズイユ流（初期流速ゼロから一定圧力勾配で定常状態へ遷移）。

\textbf{合否判定:} 定常状態の流速プロファイルが放物線（解析解）と完全に一致すること。

\subsection*{3.3 圧力ポアソン方程式（PPE）ソルバのテスト（関連: 9.2節）}

\textbf{検証対象:} 行列生成関数および連立一次方程式ソルバ（BiCGSTABやGMRES等）。

\textbf{検証方法:} 密度比1:1000のステップ状密度分布を持つ静止場に対し、人為的な発散（$\nabla \cdot \mathbf{u} \neq 0$）を与えて圧力を解き、速度場を修正（射影）する。

\textbf{合否判定:} 修正後の速度場の発散 $|\nabla \cdot \mathbf{u}^{n+1}|$ が、領域内の全セルで $10^{-10}$ 以下になること。

\section*{フェーズ4: システム統合テスト（System Validation）}

全てのモジュールを結合し、フラクショナルステップ法としてのトータル性能を検証する。

\subsection*{4.1 静止液滴問題（Static Droplet） - Balanced-Forceの究極検証}

\textbf{目的:} 「圧力勾配の離散化」と「表面張力の離散化」がセルのエッジ上で完全に釣り合っているか（演算子不整合誤差の排除）の最終確認。

\textbf{検証方法:} 無重力、初期流速ゼロの空間に静止液滴を配置し、数十タイムステップ実行する。

\textbf{合否判定:} 全セルの流速 $|\mathbf{u}|$ が、時間経過に関わらずマシンイプシロン（倍精度で $\approx 10^{-15}$）のオーダーを維持すること。これが達成できなければ、本論文の「Balanced-Force」の実装にバグがある。

\subsection*{4.2 振動液滴問題（Oscillating Droplet）}

\textbf{目的:} 時間積分（動的挙動）と粘性・表面張力のバランス確認。

\textbf{合否判定:} 振動周期がLambの理論解と一致し、振幅の対数減衰率から逆算した粘性係数が設定値と一致すること。

\subsection*{4.3 上昇気泡問題（Rising Bubble Benchmark）}

\textbf{目的:} 密度比・粘度比が存在し、重力、表面張力、粘性が複雑に絡み合う実用問題での妥当性確認。

\textbf{検証方法:} Hysing et al. (2009) のベンチマーク設定（Case 1 \& Case 2）に準拠。

\textbf{合否判定:} 気泡の終端上昇速度、重心位置の時間変化、および最終的な気泡の「円形度（Circularity）」が、ベンチマーク論文のリファレンス値と許容誤差内で一致すること。

\section*{結言}

本仕様書に定めたフェーズ1からフェーズ3までの「単体テスト」を一つ一つクリア（自動テスト化を推奨）することで、理論式からコードへの落とし込みにおけるバグを完全に排除できる。これらをクリアした上ではじめて、フェーズ4の物理現象シミュレーションに移行すべきである。統合後に寄生流れや質量欠損等の不具合が生じた場合でも、本テスト体系に立ち戻ることで、原因の切り分けが迅速に実行可能となる。

\end{document}